\documentclass[UTF8,a4paper,11pt]{ctexart}
\usepackage[margin=25mm,headheight=15pt]{geometry}
\usepackage{amsmath,amssymb,amsthm,fancyhdr,hyperref,xcolor}
\hypersetup{colorlinks=true,urlcolor=blue!55!black,linkcolor=black}
\setlength{\parindent}{2em}
\setlength{\parskip}{0.25em}
\linespread{1.13}
\pagestyle{fancy}
\fancyhf{}
\fancyhead[L]{\small 三千弱水 · 高中导数真题}
\fancyhead[R]{\small 2024 年 · 首轮独立复核}
\fancyfoot[C]{\thepage}
\renewcommand{\headrulewidth}{0.3pt}
\newcommand{\qtitle}[1]{\par\bigskip\noindent{\large\bfseries #1}\par\medskip}
\newcommand{\sourceurl}[1]{\par\noindent{\small\color{black!65}来源：\href{#1}{固定版本 LaTeX 源码（点击查看）}}\par}
\begin{document}
\begin{center}
  {\LARGE\bfseries 高考导数解答题：独立核验选编}\par\vspace{0.6em}
  {\large 2024 年新高考两道解答题}\par\vspace{0.8em}
  {\small 仅收录数学表达和独立推导；并非原试卷复刻}\par
\end{center}
\vspace{0.5em}
\noindent\textbf{审核说明：}已对下列两题给出的结论、关键代数推导、参数端点和定义域进行独立核查；未声称原试卷全文、上游仓库其他题目均正确。原题以所列源码链接为准。\par
\qtitle{一、2024 年新高考一卷第 18 题}
设
\[
 f(x)=\ln\frac{x}{2-x}+ax+b(x-1)^3,\qquad 0<x<2,
\]
其中 $a,b\in\mathbb R$。分别研究：(1) $b=0$ 且导数处处非负时 $a$ 的最小值；(2) 曲线的中心对称性；(3) $f(x)>-2$ 恰在 $1<x<2$ 成立时 $b$ 的范围。

\noindent\textbf{解答与核验。}

\noindent\textbf{(1)} 当 $b=0$ 时，
\[
 f'(x)=\frac{1}{x}+\frac{1}{2-x}+a=\frac{2}{x(2-x)}+a.
\]
由于 $x(2-x)=1-(x-1)^2\le 1$ 且 $x(2-x)>0$，有 $\frac{2}{x(2-x)}\ge 2$，等号在 $x=1$ 取到。因此 $f'(x)\ge0$ 对所有 $x\in(0,2)$ 成立当且仅当 $a\ge-2$，故 $a_{\min}=-2$。

\noindent\textbf{(2)} 注意到定义域 $(0,2)$ 在变换 $x\mapsto2-x$ 下不变，且
\[
 f(2-x)=-\ln\frac{x}{2-x}+a(2-x)-b(x-1)^3=2a-f(x).
\]
故 $\bigl(x,f(x)\bigr)$ 与 $\bigl(2-x,2a-f(x)\bigr)$ 均在曲线上，二者的中点为 $(1,a)$，曲线关于 $(1,a)$ 中心对称。

\noindent\textbf{(3)} 函数在 $x=1$ 连续，且题设要求 $x\le1$ 时 $f(x)\le-2$、$x>1$ 时 $f(x)>-2$。左右逼近 $1$ 可得 $f(1)=-2$，即 $a=-2$。记 $t=x-1\in(-1,1)$，则
\[
 f(1+t)+2=H_b(t):=\ln\frac{1+t}{1-t}-2t+bt^3.
\]
函数 $H_b$ 为奇函数，因而题设等价于 $H_b(t)>0$ 对一切 $0<t<1$ 成立。令
\[
 R(t)=\ln\frac{1+t}{1-t}-2t-\frac23t^3.
\]
对 $0<t<1$，
\[
 R'(t)=\frac{2}{1-t^2}-2-2t^2
       =\frac{2t^4}{1-t^2}>0,
 \qquad R(0)=0.
\]
于是 $R(t)>0$，且
\[
 H_b(t)=R(t)+\left(b+\frac23\right)t^3.
\]
若 $b\ge-\frac23$，则 $H_b(t)>0$。反之，
\[
 \lim_{t\to0^+}\frac{H_b(t)}{t^3}=b+\frac23,
\]
若 $b<-\frac23$，则充分小的正 $t$ 使 $H_b(t)<0$，矛盾。故 $\boxed{b\ge-\frac23}$。
\sourceurl{https://github.com/DxAThing/Gaokao-Math-Problems-Compilation/blob/f1ae169b58a75b3db6e35c670b3f0d5c78718188/content/2024/new_gaokao_paper_1.tex}

\qtitle{二、2024 年新高考二卷第 16 题}
设 $f(x)=e^x-ax-a^3$。分别求：(1) $a=1$ 时曲线在横坐标 $1$ 处的切线；(2) 存在严格负的极小值时参数 $a$ 的范围。

\noindent\textbf{解答与核验。}

\noindent\textbf{(1)} $a=1$ 时，$f'(x)=e^x-1$，$f(1)=e-2$，$f'(1)=e-1$。故切线为
\[
 y-(e-2)=(e-1)(x-1),\qquad\boxed{y=(e-1)x-1}.
\]
\noindent\textbf{(2)} 由 $f'(x)=e^x-a$、$f''(x)=e^x>0$，当 $a\le0$ 时 $f'(x)>0$，不存在极小值；当 $a>0$ 时唯一极小值点为 $x=\ln a$，极小值为
\[
 f(\ln a)=a-a\ln a-a^3=a\bigl(1-\ln a-a^2\bigr).
\]
因 $a>0$，极小值严格小于 $0$ 当且仅当 $G(a)=a^2+\ln a-1>0$。又
\[
 G'(a)=2a+\frac1a>0\quad(a>0),\qquad G(1)=0,
\]
所以所求为 $\boxed{a\in(1,+\infty)}$。
\sourceurl{https://github.com/DxAThing/Gaokao-Math-Problems-Compilation/blob/f1ae169b58a75b3db6e35c670b3f0d5c78718188/content/2024/new_gaokao_paper_2.tex}

\bigskip\noindent{\small\color{black!65}资料溯源：上游项目 DxAThing/Gaokao-Math-Problems-Compilation；其原创编辑与编排部分适用 CC BY-SA 4.0，部分原试题等第三方材料可能另有权利。本文推导为重新独立撰写，原项目与作者署名应继续保留。}
\end{document}